\chapter{Layer 6: Logic Validation and Static Analysis}
\label{ch:static}

\section{Purpose and Contract}
\label{sec:static-purpose}

Layer~6 is the first \emph{defensive} boundary after computation: it checks
that the logic is consistent and the code is safe. Its inbound contract is
``a set of obligations, a kernel, and a codebase''; its outbound contract is
``a consistency verdict and a static-analysis report.'' Two tools implement
it: a first-order logic resolution checker and a CodeQL static-analysis
pass.

\section{FOL Resolution Checker}
\label{sec:static-fol}

The \texttt{layers/fol/} module is a classical-logic theorem prover. Its
\texttt{README} states the contract plainly:
\begin{quote}
\textbf{Layer 6: FOL Resolution Checker} \\
Purpose: Verified classical logic theorem prover. \\
Status: 15/15 theorems verified.
\end{quote}
Resolution is the refutation-complete inference rule: from clauses
$(A \lor l)$ and $(\neg l \lor B)$ infer $(A \lor B)$. A goal is proven by
negating it, converting to clause form, and showing that resolution yields the
empty clause. Because resolution is sound and refutation-complete for
first-order logic, a verified implementation gives Layer~6 the same
trustworthy status as the Lean obligations upstream --- independent confirmation
from a different proof method.

\section{Static Analysis with CodeQL}
\label{sec:static-codeql}

The \texttt{layers/codeql/} module applies static analysis to the
implementation. Where Lean reasons about mathematics, CodeQL reasons about
the program text: it queries the abstract syntax graph for dangerous patterns
--- unchecked casts, missing bounds checks, tainted inputs reaching the kernel.
The presence of a dedicated static-analysis layer means a class of
implementation bugs is caught at the boundary rather than in production.

\section{Why Both}
\label{sec:static-why}

The two tools answer different questions:
\begin{itemize}
  \item FOL: \emph{Is the logic internally consistent?} (mathematical)
  \item CodeQL: \emph{Is the implementation free of known defect classes?}
        (engineering)
\end{itemize}
Running both realizes defense in depth: a proof can be logically correct yet
implemented unsafely, and vice versa. The layer fails unless \emph{both} pass.

\section{Integration with the Pipeline}
\label{sec:static-int}

In the CI scaffold (\cref{ch:build}), Layer~6 corresponds to the
``Logic validation and static analysis'' step. Its reports are sealed as
receipts (Layer~9) and consulted by the policy layer (Layer~7) before any
artifact is permitted to publish.

\section{Invariant Inference}
\label{sec:static-infer}

The static pass does not require hand-written invariants; it infers them
from the C$--$ control flow. For the loop in \cref{lst:work-cmm} the
inferred invariant
\[
0 \le i0 \le 4 \;\wedge\; r0 = \sum_{k=0}^{i0-1} v_k
\]
is derived by a standard abstract interpretation over integer intervals,
then strengthened by the FOL layer (\cref{sec:work-fol}) into a
prover-ready clause. The same machinery scales to the matrix
contraction, where the invariant tracks the contracted row index.
